\begin{annexitem}{annex:p7.real.1}{thm:p07-number-realizability-soundness}{Uniform number realizers}
\begin{proof}
Use the acceptable-numbering package of \cref{def:p04-ce-numbering},
its universal evaluator, the effective parameterization proved in
\cref{thm:p04-local-parameter}, and the fixed primitive-recursive
pairing and projections. Finite parameter tuples are coded with that
pairing. If \(e\) is an index, write \(\{e\}(a)\) for its
partial value at \(a\). A displayed program depending on numerical
parameters can be compiled effectively to an index with those parameters
fixed. Programs may be undefined on inputs that do not realize their
stated premises. Every termination assertion below concerns an input
satisfying that premise.

We prove the stronger, uniform statement: for a derivation of
\(A(\vec n)\), the construction supplies a recursive procedure
\(r_A(\vec n)\) which terminates for every numerical parameter tuple
and returns a realizer of its numeral instance. It suffices to verify
the twelve groups of logical rules used in the preceding Dialectica
proof, equality, and the handbook's arithmetic axioms. Parameterization
is performed at compilation time; it does not require evaluating a
program at every possible future argument. Intermediate premises may
have extra numerical parameters that do not occur in the conclusion.
After the rule's required eigenvariable abstraction or witness
substitution, supply zero for any such remaining parameter. Uniformity
of the child procedures makes this specialization legitimate. An empty
parameter tuple can be represented by a fixed dummy numerical input.

\emph{1. Modus ponens.} Compute the realizing index
\(f=r_{A\to B}(\vec n)\) and \(a=r_A(\vec n)\), and return
\(\{f\}(a)\). By the implication clause this computation terminates
and realizes \(B(\vec n)\).

\emph{2. Composition.} Given realizing indices \(f\) for
\(A\to B\) and \(g\) for \(B\to C\), compile the program
\(a\mapsto\{g\}(\{f\}(a))\). On an \(A\)-realizer its
inner computation terminates with a \(B\)-realizer, and the outer
one then terminates with a \(C\)-realizer.

\emph{3. Contractions.} The program \(a\mapsto\langle a,a\rangle\)
realizes \(A\to A\land A\). The program
\(e\mapsto\pi_1(e)\) realizes \(A\lor A\to A\), since either
valid tag carries an \(A\)-realizer. Unlike the negative-challenge
interpretation, this operation needs no truth test of \(A\).

\emph{4. Injection and projection.} Use
\(a\mapsto\langle0,a\rangle\) for \(A\to A\lor B\), and
\(e\mapsto\pi_0(e)\) for \(A\land B\to A\).

\emph{5. Commutations.} Exchange the components of a conjunction.
For a disjunction retain its payload and change its tag from zero to
one or from one to zero. Its realizing inputs have exactly these tags.

\emph{6. Disjunction monotonicity.} Given an index \(f\) realizing
\(A\to B\), a program for \(C\lor A\to C\lor B\) leaves
a zero-tagged input unchanged and sends a one-tagged payload \(a\)
to \(\langle1,\{f\}(a)\rangle\). The only new evaluation is
on an \(A\)-realizer and hence terminates correctly.

\emph{7. Currying.} If \(f\) realizes \(A\to(B\to C)\), send
\(\langle a,b\rangle\) to \(\{\{f\}(a)\}(b)\). If \(f\)
realizes \(A\land B\to C\), send \(a\) to an index for the
program \(b\mapsto\{f\}(\langle a,b\rangle)\). Effective
parameterization supplies that index. Both directions satisfy the
implication clauses exactly.

\emph{8. Falsity.} No number realizes the false atomic equation
\(0=1\). Therefore any fixed program index realizes
\(\bottom\to C\), whose condition on input realizers is vacuous.

\emph{9. Universal introduction.} Let \(f(\vec n,z)\) realize
\(A\to B(z)\), with \(z\) absent from \(A\). Given an
\(A\)-realizer \(a\), return an index for
\[
 z\longmapsto\{f(\vec n,z)\}(a).
\]
Each evaluation terminates by the induction hypothesis and realizes
\(B(z)\). The returned index therefore realizes \(\forall zB(z)\).

\emph{10. Universal elimination.} For
\(\forall zA(z)\to A(t)\), evaluate the input index at the
numerical value of the arithmetic term \(t\). The universal
realizability clause gives termination and the desired numeral instance.

\emph{11. Existential introduction.} For
\(A(t)\to\exists zA(z)\), pair the numerical value of \(t\)
with the input realizer.

\emph{12. Existential elimination.} Suppose
\(f(\vec n,z)\) realizes \(A(z)\to B\), with \(z\) not free
in \(B\). On input \(\langle z,a\rangle\) for
\(\exists zA(z)\), return \(\{f(\vec n,z)\}(a)\).
The witness \(z\) is supplied by the input; the induction hypothesis
and implication clause give the unchanged \(B\)-realizer.

\emph{Equality and numerical axioms.} An atomic formula is realized
by every number when it is true, and by no number when it is false.
Reflexivity is therefore realized by zero. A program returning a
constant program returning zero realizes atomic equality replacement:
if its two successive inputs realize the equality and the replaced
atomic formula, numerical equality makes the resulting atomic formula
true. Successor's nonzero value, successor injectivity, and the defining
equations of addition and multiplication are true numerical assertions;
their atomic and implication clauses are realized by the corresponding
constant programs. For
\(n\ne0\to\exists m\,n=Sm\), ignore the input realizer and
return \(\langle p(n),0\rangle\). Such an input can exist only
when \(n\ne0\), and then \(n=Sp(n)\). These constructions are
uniform in numerical parameters, and universal closures are compiled
by the quantifier constructions above.

\emph{Induction.} An input realizer for the antecedent of
\[
 A(0)\land\forall i(A(i)\to A(Si))\to\forall iA(i)
\]
is a pair \(\langle a,f\rangle\), where \(a\) realizes \(A(0)\)
and \(f\) realizes the universal step assertion. Return an index
for the following program on input \(n\): start with \(r=a\),
and for \(i=0,\ldots,n-1\) evaluate
\(h=\{f\}(i)\) and then replace \(r\) by \(\{h\}(r)\).
Ordinary induction on this finite loop proves that it terminates and
that after \(i\) iterations \(r\) realizes \(A(i)\). Indeed the
universal clause supplies a terminating \(A(i)\to A(Si)\)-realizer,
and the implication clause then supplies its next value. Thus the
compiled program realizes \(\forall iA(i)\). Its index is obtained
effectively from \(a,f\); the formula itself is not queried as an
oracle.

All cases preserve the stated uniformity in free numerical parameters.
Induction on the finite derivation gives the recursive procedure
\(r_A\). For a closed conclusion, run it on the empty parameter
tuple and output its value. The termination argument just proved shows
that this algorithm halts on every valid HA derivation. It can be made
total on arbitrary input strings by first checking the finite source
proof and returning a default value for an invalid one. The theorem
claims effectiveness, not a primitive-recursive bound on the running
time of evaluating the extracted programs. It also makes no inference
from a realizer of an arbitrary formula to an HA proof of that formula.
\end{proof}
\end{annexitem}

\begin{annexitem}{annex:p7.real.2}{thm:p07-markov-number-realizability}{Realizing Markov's principle}
\begin{proof}
Fix the numerical parameters of the displayed predicate, and abbreviate
\(E=\exists nA(n)\) and
\(D=\forall n(A(n)\lor\neg A(n))\).
Given a realizer \(d\) of \(D\), compile a program which takes a
realizer \(e\) of \(\neg\neg E\) and searches through
\(n=0,1,2,\ldots\), evaluating \(\{d\}(n)\) at each stage.
When its tag is zero, with payload \(a\), return
\(\langle n,a\rangle\). Each individual evaluation terminates,
since \(d\) realizes the universal assertion, and a zero-tagged
payload realizes \(A(n)\).

To prove termination of the search, work in the stated classical
metatheory. If \(E\) had no realizer, every index would realize
\(\neg E\), since its implication clause would have no realizing
input. But \(e\) would then have to send such an index to a realizer
of \(\bottom\), which is impossible. Hence some
\(\langle n,a\rangle\) realizes \(E\). At that \(n\), the
value of \(\{d\}(n)\) cannot have tag one: its payload would
realize \(\neg A(n)\) and would have to send the known
\(A(n)\)-realizer \(a\) to a realizer of \(\bottom\).
Thus a zero tag occurs, and the sequential search eventually finds one.
The returned pair realizes \(E\).

The program does not inspect \(e\); the promise that \(e\) realizes
the double negation is what guarantees termination. This is legitimate
for realizability of an implication, which requires correct termination
only on realizing inputs. Effective parameterization supplies the
outer index taking \(d\) to this program, uniformly in any numerical
parameters. It realizes exactly \(D\to(\neg\neg E\to E)\).
The argument proves validation by this semantics, not derivability of
the scheme in HA and not an independence result.
\end{proof}
\end{annexitem}

\begin{annexitem}{annex:p7.real.3}{thm:p07-ipc-disjunction-property}{Propositional disjunction extraction}
\begin{proof}
Apply \cref{thm:p07-lj-exact-cut-elimination} to the given
propositional derivation of \(\Rightarrow A\lor B\). Its scope
excludes both Cut and Mix, and its conclusion preserves the same
literal endsequent. Applied to a propositional proof, its construction
introduces no quantifier: the multicut reductions only decompose the
formulas encountered in that proof. We therefore obtain a cut-free
propositional proof in the same sheet. By
\cref{thm:p07-lj-disjunction-last-rule}, its last inference supplies
an empty-antecedent proof of \(A\) or of \(B\). No arithmetic
axioms or equality rules are involved in this argument.
\end{proof}
\end{annexitem}

\begin{annexitem}{annex:p7.real.4}{thm:p07-proof-carrying-soundness}{Constructing the carried proofs}
\begin{proof}
We enrich the uniform program construction of
\cref{thm:p07-number-realizability-soundness}. A realizer is now a
pair of a proof code and a payload. Whenever a program examines logical
structure it examines the payload, and every recursive input and output
is an entire proof-carrying realizer. It is not enough to attach the
original proof at the outermost node; the constructions below produce
the component proofs as well.

\emph{Finite proof operations.} Modus ponens, substitution of closed
numerals for free numerical parameters, and the fixed logical axiom
instances are effective operations on HA proof codes. Eigenvariables
are renamed first to preserve their side conditions. In particular
there are finite proof constructors for conjunction introduction and
projection, disjunction introduction and commutation, existential
introduction with a specified numeral, and universal elimination at
a specified numeral. For clarity, conjunction introduction requires no
extra logical axiom: derive
\((A\land B)\to(A\land B)\) from duplication, projection, and
composition; the currying rule gives \(A\to(B\to A\land B)\).
Two applications of modus ponens then combine proofs of \(A\) and
\(B\). The other listed operations use the corresponding displayed
logical axiom and modus ponens. Thus these are literal finite proof
constructions in the same HA presentation.

There is also an effective proof of every true closed numerical
equality. Evaluate its terms and construct proofs that each term equals
its computed numeral: induction on the term uses the defining successor,
addition and multiplication equations, with finitely many numeral
recursion steps. Equal computed numerals are then joined by equality
replacement and transitivity. This operation is used only when the
calculation has verified the closed equality. It is not a procedure
turning arbitrary arithmetic truth into provability.

\emph{Invariant.} At a proof node \(D:A(\vec z)\), the desired
procedure returns a proof-carrying realizer of every numeral instance.
The current derivation, after numeral substitution, supplies a finite
proof \(p_D\) of that whole instance. When the construction forms
a new outer realizer it can use this proof as its first component.
Extra numerical parameters of intermediate premises are specialized to
zero after the eigenvariable operations, as in ordinary soundness.
Proof substitution is the capture-free operation with freshened
eigenvariables described above. The recursively required payload
components are constructed as follows.

\emph{Application and composition.} On a proof-carrying input, a
realizing implication's program already returns a proof-carrying result.
Thus modus ponens and composition preserve the invariant with no new
assumption about provability. Their source proof nodes also supply the
corresponding whole-formula proofs. The program indices are compiled
by effective parameterization, just as before.

\emph{Conjunctions and disjunctions.} Pairing two realizers is now
accompanied by the conjunction proof constructed from their two carried
proofs. A projection returns the appropriate whole subrealizer.
Disjunction introduction attaches a proof of the disjunction obtained
from the selected component's carried proof, and keeps that component
as its payload. Contraction of a disjunction returns its selected
subrealizer, which already carries the necessary disjunct proof;
duplication of a conjunction uses its input proof twice. Commutation
rearranges existing components and obtains the new outer proof by its
fixed logical theorem. For disjunction monotonicity, leave the retained
branch's subrealizer intact, or run the supplied implication on the
other branch; disjunction introduction then supplies the appropriate
new outer proof. These operations verify logical groups 3--6 without
extracting a disjunct from an arbitrary proof alone.

\emph{Currying.} To uncurry, take the two subrealizers of a conjunction
and apply the first realizing implication to one, then its returned
implication to the other. To curry a realizer for \(A\land B\to C\),
take a proof-carrying \(A\)-realizer \(a\) and compile the program
which, on a proof-carrying \(B\)-realizer \(b\), builds their
proof-carrying conjunction and applies the given program. The proof
of \(B\to C\) attached to this compiled program is obtained by
modus ponens from the current theorem
\(A\to(B\to C)\) and the proof carried by \(a\). Thus the
returned implication is proof-carrying, rather than just an index
with the right behavior. This verifies both directions of group 7.

\emph{Falsity and equality.} No proof-carrying realizer of falsity
exists because its atomic payload condition is false. Falsity
elimination therefore has no realizing input to process. Its own
implication proof is the logical axiom. For atomic equality replacement,
the input proof of the equality and the later input proof of the
replaced atomic formula give a proof of the output atom by two
applications of the equality axiom and modus ponens. Numerical equality
preserves its truth, so zero is a valid atomic payload. The intermediate
implication gets its proof after the first application of modus ponens.
Reflexivity is its own carried proof with an atomic payload. This
covers group 8 and the equality axioms.

\emph{Universal introduction.} At a source inference from
\(A\to B(z)\) to \(A\to\forall zB(z)\), take a proof-carrying
\(A\)-realizer \(a\). The current theorem's proof, combined with
the proof carried by \(a\), gives one finite HA proof of
\(\forall zB(z)\). Attach it to the compiled program which, on
\(n\), obtains the enriched realizer of the child theorem
\(A\to B(\overline n)\) and applies it to \(a\). The induction
hypothesis makes every such computation terminate with a
proof-carrying \(B(\overline n)\)-realizer. The condition that
\(z\) is not free in \(A\) makes the same input \(a\) admissible
at every \(n\). The universal proof was obtained by a finite
inference and modus ponens, not by combining infinitely many proofs.

\emph{The other quantifier rules.} A universal realizer evaluated at
a numeral already returns a proof-carrying instance. Existential
introduction pairs the chosen numeral with the input instance realizer,
and attaches the existential proof obtained from its carried instance
proof. Existential elimination reads that numeral and that full
subrealizer and applies the uniformly compiled child implication.
The returned realizer carries a proof of the unchanged conclusion.
This verifies groups 9--12 with their original eigenvariable conditions.

\emph{Numerical axioms.} Closed true equations have the proof and
atomic payload described above. For successor injectivity, proof
application supplies the output equality proof and its numerical
truth follows from injectivity. Successor nonzero has a vacuous
negative payload and its axiom proof. For predecessor existence,
a realizing input of \(n\ne0\) is impossible at \(n=0\): the
true equation \(0=0\) has a proof-carrying realizer, which such an
input would have to send to falsity. At a positive numeral, compute
its predecessor and the finite equational proof of
\(\overline n=S\overline{p(n)}\). This gives a proof-carrying
instance and hence a proof-carrying existential by the preceding
constructor. The implication and its universal closures carry their
own source axiom proofs.

\emph{Induction.} On a proof-carrying realizer of
\(A(0)\land\forall i(A(i)\to A(Si))\), run the same finite
iteration program as in the ordinary soundness proof, but on whole
proof-carrying realizers. Every step receives a carried proof of
\(A(i)\) and returns one of \(A(Si)\), with the required payload,
by the induction hypothesis on the loop length. To supply the proof
component for the output \(\forall iA(i)\), apply the current
induction-axiom proof to the proof carried by the input conjunction.
This is a single finite modus ponens inference. Attach that proof to
the compiled iteration program. The root realizer of the induction
axiom carries the induction axiom's own finite proof. No arithmetic
normalization or infinitary inference has entered the construction.

This accounts for every source proof node and every newly formed
compound payload in the twelve logical groups, equality, and the
arithmetic axioms. Induction on the source derivation therefore proves
the uniform invariant. All proof manipulation is finite and effective;
termination of the evaluated programs follows on the stated realizing
inputs just as in ordinary soundness. The construction never assumes
that a source proof already supplies the disjunct or numerical witness
being sought. Those are produced by the recursively constructed
payload, whose component proofs have now been verified.
\end{proof}
\end{annexitem}

\begin{annexitem}{annex:p7.real.5}{thm:p07-ha-disjunction-property}{Extracting an HA disjunct proof}
\begin{proof}
Apply \cref{thm:p07-proof-carrying-soundness} to the given closed
HA proof of \(A\lor B\). The resulting realizer has a payload
with tag zero or one and a proof-carrying realizer of the selected
disjunct. Read the proof component of that subrealizer. It is an HA
proof of \(A\) or of \(B\), respectively. This is not the
unsupported inference from an ordinary number realizer to a proof;
the preceding theorem constructed and verified this particular
proof-carrying payload.
\end{proof}
\end{annexitem}

\begin{annexitem}{annex:p7.real.6}{thm:p07-ha-numerical-existence}{Extracting a numeral and its proof}
\begin{proof}
Apply \cref{thm:p07-proof-carrying-soundness} to the closed proof of
\(\exists nA(n)\). Its existential payload contains a standard
number \(k\) and a proof-carrying realizer of \(A(\overline k)\).
Read the proof component of that subrealizer. It is the required finite
HA derivation of the numeral instance. Computing the realizer, reading
\(k\), and reading its carried proof give an effective extraction
procedure on valid source proofs; no bound on its running time is
asserted.
\end{proof}
\end{annexitem}
